% ===================== 写作 Writing =====================
\ptesection{写作}{Writing · Summarize Written Text + Write Email}{71}

\vspace{1mm}
{\small\color{gray} 本部分含 \textbf{SWT 概括写作} 2 题（满分各 8）与 \textbf{WE 邮件写作} 2 题（满分各 15）。
SWT 仅展示“你的作答 + 批改”；WE 另含平台“范文”。下方每题保留逐项评分与 AI 中文建议，并为原文 / 题干 / 范文附\textbf{中文翻译}。本套写作总分 \textbf{71 / 90}。}
\vspace{2mm}

\newcolumntype{Y}{>{\RaggedRight\arraybackslash}X}

% ---------------- Q1 ----------------
\begin{qblock}{Q1\quad SWT \#60\ \ Elizabeth Blackwell}{8 / 8\quad·\quad 用时 05:28}

\ptesub{原文 Source Text}
\begin{promptbox}
Elizabeth Blackwell's journey to becoming the first woman to receive a medical degree in the United States is a story of resilience and determination in the face of immense obstacles. Born in 1821 in Bristol, England, Blackwell moved to the United States with her family at a young age. Her interest in medicine was sparked by a dying friend who expressed that a female physician would have made her medical experience less traumatic.

Despite widespread prejudice against women in medicine at the time, Blackwell was undeterred. She applied to several medical schools, facing rejection after rejection. Finally, she was admitted to Geneva Medical College in New York, where she graduated first in her class in 1849, breaking barriers for women in the field of medicine.

Blackwell's contributions extended beyond her personal achievements. She was an ardent advocate for women in medicine, establishing the New York Infirmary for Women and Children with her sister, Dr.\ Emily Blackwell, and Dr.\ Marie Zakrzewska. This institution not only provided medical care to the underserved but also offered training and experience for women medical practitioners.
\end{promptbox}
\zh{伊丽莎白·布莱克威尔成为美国第一位获得医学学位的女性，她的历程是一个在巨大阻碍面前坚韧不拔、矢志不渝的故事。她于 1821 年生于英格兰布里斯托尔，年幼时随家人移居美国。她对医学的兴趣，源于一位临终好友——这位朋友曾说，若当年有女医生为她诊治，她的就医经历或许不会那么痛苦。尽管当时社会对女性从医普遍存有偏见，布莱克威尔却毫不退缩。她申请了多所医学院，屡遭拒绝，最终被纽约的日内瓦医学院录取，并于 1849 年以全班第一名的成绩毕业，为女性在医学领域打破了壁垒。布莱克威尔的贡献远不止于个人成就。她是女性从医的热忱倡导者，与妹妹埃米莉·布莱克威尔医生及玛丽·扎克热夫斯卡医生一同创办了纽约妇女儿童医院；这所机构不仅为弱势群体提供医疗服务，也为女性医务工作者提供了培训与实践的机会。}

\ptesub{你的作答 Your Response}
\begin{youbox}
Elizabeth Blackwell's journey is a story of resilience and determination in the face of immense obstacles, and Blackwell's contributions extended beyond her personal achievements, and she was an ardent advocate for women in medicine.
\end{youbox}

\ptesub{评分明细 Score Breakdown}
{\small\scorehead
\begin{tabularx}{\linewidth}{l c Y}
\rowcolor{cHead}\textcolor{white}{\bfseries 维度} & \textcolor{white}{\bfseries 得分} & \textcolor{white}{\bfseries AI 建议}\\
Content 内容        & \sfull{2 / 2} & 内容很棒\\
Form 格式          & \sfull{2 / 2} & Form 很棒，这是必拿分数\\
Grammar 语法       & \sfull{2 / 2} & 很棒，AI 小猩几乎检查不出什么语法问题\\
Vocabulary 词汇    & \sfull{2 / 2} & 很棒，再接再厉\\
\end{tabularx}}
\smallskip
\textbf{本题得分：}\ {\color{cAccent}\bfseries 8 / 8}\quad{\footnotesize\color{gray}（满分，无批改）}

\end{qblock}

% ---------------- Q2 ----------------
\begin{qblock}{Q2\quad SWT \#69\ \ Beta Testing}{8 / 8\quad·\quad 用时 03:36}

\ptesub{原文 Source Text}
\begin{promptbox}
There are over a million apps in the world. Most of these apps have to go through mobile beta testing for them to be relevant in the market. Beta testing allows testing your app's real-world performance. This allows quality assurance teams to identify and focus on any bugs that will appear. By doing so, it's assumed that customers can expect the product to work without any problems.

Understanding the problems found ahead of time makes the development team efficient. They improved the quality of the product before release. This means they can improve on it during the development process to ensure better quality products at launch. Running a beta test on your product decreases the time for it to market. Understanding customer expectations improve other aspects of the product and improve other processes.
\end{promptbox}
\zh{世界上有超过一百万款应用程序。其中大多数都必须经过移动端的 Beta 测试，才能在市场上保持竞争力。Beta 测试可以检验应用在真实环境中的表现，让质量保障团队能够发现并集中处理可能出现的各种 bug。如此一来，便可假定用户能够期待产品正常、无故障地运行。提前发现问题能让开发团队更高效；他们在发布前就提升了产品质量，这意味着可以在开发过程中不断改进，从而确保上线时交付更优质的产品。对产品进行 Beta 测试可缩短其上市时间；而了解客户的期望，则有助于改善产品的其他方面、优化其他流程。}

\ptesub{你的作答 Your Response}
\begin{youbox}
There are over a million apps in the world, and most of them have to go through mobile beta testing for them to be relevant in the market, and understanding the problems found ahead of time makes the development team efficient.
\end{youbox}

\ptesub{评分明细 Score Breakdown}
{\small\scorehead
\begin{tabularx}{\linewidth}{l c Y}
\rowcolor{cHead}\textcolor{white}{\bfseries 维度} & \textcolor{white}{\bfseries 得分} & \textcolor{white}{\bfseries AI 建议}\\
Content 内容        & \sfull{2 / 2} & 内容很棒\\
Form 格式          & \sfull{2 / 2} & Form 很棒，这是必拿分数\\
Grammar 语法       & \sfull{2 / 2} & 很棒，AI 小猩几乎检查不出什么语法问题\\
Vocabulary 词汇    & \sfull{2 / 2} & 很棒，再接再厉\\
\end{tabularx}}
\smallskip
\textbf{本题得分：}\ {\color{cAccent}\bfseries 8 / 8}\quad{\footnotesize\color{gray}（满分，无批改）}

\end{qblock}

% ---------------- Q3 ----------------
\begin{qblock}{Q3\quad WE \#10\ \ Healthier Work Environment}{12.9 / 15\quad·\quad 用时 07:43}

\ptesub{题干 Task Prompt}
\begin{promptbox}
Your company has recently announced a new initiative to improve employee well-being. Write an email to your supervisor suggesting three ways to promote a healthier and more balanced work environment. You should write between 80 and 120 words.
\end{promptbox}
\zh{贵公司近期宣布了一项旨在提升员工福祉的新举措。请给你的主管写一封邮件，建议三种方式来营造更健康、更平衡的工作环境。字数应在 80 至 120 之间。}

\ptesub{范文 Model Answer（平台参考）}
\begin{modelbox}
Dear Mrs.\ Rivera,

I hope this email finds you well. With the recent announcement of the new initiative to improve employee well-being, I would like to suggest the following ways to promote a healthier and more balanced work environment:

Firstly, I recommend implementing flexible work hours to allow employees to better manage their work-life balance and reduce stress.

Secondly, wellness programs such as yoga classes, meditation sessions, or health workshops can be introduced to support employees' physical and mental well-being.

Finally, providing healthy snack options in the office can encourage better eating habits and contribute to overall employee wellness.

Thank you for considering these suggestions.

Best regards,\\ Peter Pan
\end{modelbox}
\zh{尊敬的 Rivera 女士：您好。随着公司近期宣布提升员工福祉的新举措，我想就营造更健康、更平衡的工作环境提出以下几点建议。第一，我建议推行弹性工作时间，让员工更好地平衡工作与生活、缓解压力。第二，可以引入瑜伽课、冥想课或健康讲座等健康项目，以支持员工的身心健康。最后，在办公室提供健康零食，可鼓励更好的饮食习惯，促进员工整体的健康。感谢您考虑这些建议。此致敬礼，Peter Pan}

\ptesub{你的作答 Your Response（含批改）}
\begin{youbox}
Dear \fix{supervisor}{Supervisor},

I hope this email finds you well. \fix{I'mwriting}{I'm writing（缺空格）} to suggest three ways to promote a healthier and more balanced work environment.

Firstly, we can go to work earlier so that we can go home earlier.

Secondly\fix{, company}{, the companies} can prepare more food\ins{,} so the work environment will be \del{more} healthier.

Lastly, we can hold more activities, which can improve employee well-being.

I would appreciate it if you could consider my suggestions. Looking \add{forwards}{forward}\,(固定搭配 look forward to) to your reply.

Best regards,\\ Kathy
\end{youbox}
{\footnotesize\color{gray} 注：橙色波浪线（补）为平台\textbf{漏标}、本卷补标的错误。``Looking forwards to'' 的习惯搭配为 ``look forward to''，建议用 ``forward''。}

\ptesub{评分明细 Score Breakdown}
{\small\scorehead
\begin{tabularx}{\linewidth}{l c Y}
\rowcolor{cHead}\textcolor{white}{\bfseries 维度} & \textcolor{white}{\bfseries 得分} & \textcolor{white}{\bfseries AI 建议}\\
Content 内容              & \sfull{3 / 3}   & 内容很棒\\
Form 格式                & \sfull{2 / 2}   & Form 很棒，这是必拿分数\\
Email Conventions 邮件格式 & \sfull{2 / 2}   & 邮件书写格式很棒！\\
Organization 结构         & \sfull{2 / 2}   & 段落分明，文章架构清晰明了，很棒\\
Vocabulary 词汇           & \spart{1.5 / 2} & PTE 对词汇不要求很难，但需要恰当且拼写正确。点击彩色单词查看错误和修改建议\\
Grammar 语法              & \szero{0.4 / 2} & 语法在 PTE 写作中占分比很重，但其实又非常容易提高。丢分这么严重，有点不应该哦。建议找老师咨询下迅速提高的技巧\\
Spelling 拼写             & \sfull{2 / 2}   & 虽然你拼写拿了满分，但是你还是犯了 1 个拼写错误。下次注意避免哦\\
\end{tabularx}}
\smallskip
\textbf{本题得分：}\ {\color{cAccent}\bfseries 12.9 / 15}\quad{\footnotesize\color{gray}（平台标注：supervisor$\to$Supervisor 首字母大写、I'm writing 缺空格、, company$\to$, the companies、food 后缺逗号、more 应删；\textcolor{cAdd}{补标：Looking forwards$\to$forward}）}

\end{qblock}

% ---------------- Q4 ----------------
\begin{qblock}{Q4\quad WE \#11\ \ Remote Team Collaboration and Communication}{12.5 / 15\quad·\quad 用时 07:43}

\ptesub{题干 Task Prompt}
\begin{promptbox}
Your company is looking to improve team collaboration and communication in remote work settings. Write an email to your team leader, suggesting three strategies to foster better teamwork and enhance communication among remote team members. You should write between 80 and 120 words.

\smallskip Your suggestions should focus on the following three themes:
\begin{itemize}[leftmargin=1.4em,topsep=2pt,itemsep=0pt]
  \item Implementing a weekly virtual team meeting;
  \item Adopting a collaborative online platform;
  \item Scheduling regular one-on-one check-ins.
\end{itemize}
You should include all three themes. Provide supporting ideas for your suggestions.
\end{promptbox}
\zh{贵公司希望改善远程办公环境下的团队协作与沟通。请给你的团队负责人写一封邮件，建议三种策略来促进更好的团队合作、加强远程团队成员之间的沟通。字数应在 80 至 120 之间。你的建议应围绕以下三个主题：推行每周一次的线上团队会议；采用协作式在线平台；安排定期的一对一沟通。三个主题都要涵盖，并为你的建议提供支撑性想法。}

\ptesub{范文 Model Answer（平台参考）}
\begin{modelbox}
Dear Mr Harrison,

I hope this email finds you well. As we navigate the challenges of remote work, I believe we can enhance our team's collaboration and communication with a few strategic initiatives.

Firstly, implementing a weekly virtual team meeting could help us stay connected and align on projects. Secondly, adopting a collaborative online platform would streamline our workflow and ensure transparency. Lastly, scheduling regular one-on-one check-ins could support individual team members' needs and foster a stronger team bond.

I'm confident these strategies will significantly improve our remote work dynamic.

Looking forward to your feedback.

Best,\\ Peter Pan
\end{modelbox}
\zh{尊敬的 Harrison 先生：您好。在我们应对远程办公种种挑战之际，我相信通过几项有针对性的举措，可以增强团队的协作与沟通。第一，推行每周一次的线上团队会议，有助于我们保持联系、在各项目上达成一致。第二，采用协作式在线平台，能够简化工作流程、确保信息透明。最后，安排定期的一对一沟通，可以照顾到每位团队成员的需求，增进更紧密的团队凝聚力。我相信这些策略将显著改善我们的远程办公状态。期待你的反馈。此致，Peter Pan}

\ptesub{你的作答 Your Response（含批改）}
\begin{youbox}
Dear \fix{leader}{Leader},

I hope this email finds you well. \fix{I'mwriting}{I'm writing（缺空格）} to \fix{suggesting}{suggest} three strategies to foster better teamwork and enhance communication among remote team members.

Firstly, we can implement a weekly virtual team meeting, which can \fix{improves}{improve} team collaboration.

Secondly, we can \fix{adopting}{adopt} a collaborative online platform so that everyone can share their suggestions on it.

Lastly, we can \fix{scheduling}{schedule} regular one-on-one check-ins, which can \fix{improves}{improve} team communication in remote work settings.

I would appreciate it if you could consider my suggestions. Looking \add{forwards}{forward}\,(固定搭配 look forward to) to your reply.

Best regards,\\ Kathy
\end{youbox}
{\footnotesize\color{gray} 注：本题失分集中在\textbf{语法（动词形式）}——三处“can / to + 动词”后误用了 -ing / -s 形式。橙色波浪线（补）同 Q3：``forwards''$\to$``forward''。}

\ptesub{评分明细 Score Breakdown}
{\small\scorehead
\begin{tabularx}{\linewidth}{l c Y}
\rowcolor{cHead}\textcolor{white}{\bfseries 维度} & \textcolor{white}{\bfseries 得分} & \textcolor{white}{\bfseries AI 建议}\\
Content 内容              & \sfull{3 / 3} & 内容很棒\\
Form 格式                & \sfull{2 / 2} & Form 很棒，这是必拿分数\\
Email Conventions 邮件格式 & \sfull{2 / 2} & 邮件书写格式很棒！\\
Organization 结构         & \sfull{2 / 2} & 段落分明，文章架构清晰明了，很棒\\
Vocabulary 词汇           & \spart{1.5 / 2} & PTE 对词汇不要求很难，但需要恰当且拼写正确。点击彩色单词查看错误和修改建议\\
Grammar 语法              & \szero{0 / 2}   & 语法在 PTE 写作中占分比很重，但其实又非常容易提高。丢分这么严重，有点不应该哦。建议找老师咨询下迅速提高的技巧\\
Spelling 拼写             & \sfull{2 / 2} & 虽然你拼写拿了满分，但是你还是犯了 1 个拼写错误。下次注意避免哦\\
\end{tabularx}}
\smallskip
\textbf{本题得分：}\ {\color{cAccent}\bfseries 12.5 / 15}\quad{\footnotesize\color{gray}（平台标注：leader$\to$Leader 首字母大写、I'm writing 缺空格、suggesting$\to$suggest、improves$\to$improve（两处）、adopting$\to$adopt、scheduling$\to$schedule；\textcolor{cAdd}{补标：Looking forwards$\to$forward}）}

\end{qblock}
